%% paper67_van_der_waerden_schur_en.tex
%% Van der Waerden's Theorem and Schur Numbers:
%% Ramsey Theory, Arithmetic Progressions, and Combinatorial Colorings
%%
%% Author: Michael Fuhrmann
%% Project: Specialist Mathematics Project, Build 145
%% Date: 2026-03-12

\documentclass[12pt]{amsart}

\usepackage{amsmath, amssymb, amsthm}
\usepackage{mathtools}
\usepackage{enumitem}
\usepackage{hyperref}
\usepackage{geometry}
\usepackage{booktabs}
\usepackage{array}
\usepackage{xcolor}

\geometry{margin=2.5cm}

%% Theorem environments
\newtheorem{theorem}{Theorem}[section]
\newtheorem{lemma}[theorem]{Lemma}
\newtheorem{corollary}[theorem]{Corollary}
\newtheorem{proposition}[theorem]{Proposition}
\newtheorem{conjecture}[theorem]{Conjecture}

\theoremstyle{definition}
\newtheorem{definition}[theorem]{Definition}
\newtheorem{example}[theorem]{Example}
\newtheorem{remark}[theorem]{Remark}

%% Operators
\DeclareMathOperator{\AP}{AP}
\newcommand{\N}{\mathbb{N}}
\newcommand{\Z}{\mathbb{Z}}
\newcommand{\Q}{\mathbb{Q}}
\newcommand{\F}{\mathbb{F}}

%% Abstract before \maketitle (amsart style)
\begin{document}

\title{Van der Waerden's Theorem and Schur Numbers:\\
Ramsey Theory, Arithmetic Progressions, and Combinatorial Colorings}

\author{Michael Fuhrmann}
\address{Specialist Mathematics Project, Build 145}
\email{specialist-maths@project.local}
\date{2026-03-12}

\begin{abstract}
This paper presents a comprehensive treatment of two central results in Ramsey theory:
van der Waerden's theorem on arithmetic progressions in colored sets and Schur's theorem
on sum-free partitions of the integers. We develop the fundamental notions of arithmetic
progressions, monochromatic structures, and combinatorial colorings, then prove both
classical theorems in full detail. We tabulate known van der Waerden numbers
$W(k;r)$ and Schur numbers $S(r)$, highlight the landmark results of Szemer\'edi (1975)
and Green--Tao (2004) as density and prime extensions, and discuss Rado's theorem
characterizing partition-regularity in full generality. Special attention is paid to
the frontier: $S(5) = 160$ (proved in 2017 by a SAT-solver-assisted proof), while
$S(6)$ remains unknown with only $536 \leq S(6) \leq 1836$ established. Connections
to ergodic theory (Furstenberg's proof of Szemer\'edi's theorem), analytic number theory,
and open problems such as the Erd\H{o}s--Tur\'an conjecture are surveyed.
\end{abstract}

\maketitle

\tableofcontents

%% ============================================================
\section{Introduction}
%% ============================================================

Ramsey theory is the branch of combinatorics that seeks to find unavoidable order within
large enough structures. Its philosophical core is Ramsey's maxim: \emph{complete disorder
is impossible}. Among the most celebrated theorems in this area are van der Waerden's
theorem on arithmetic progressions and Schur's theorem on sum-free colorings, both
predating the formal birth of Ramsey theory yet embodying its spirit completely.

\subsection*{Arithmetic Progressions and Colorings}

An \emph{arithmetic progression} (AP) is a sequence of integers with constant difference:
$a, a+d, a+2d, \ldots, a+(k-1)d$. A coloring of a set assigns each element one of
finitely many colors. A monochromatic AP is one whose elements all share the same color.
The central question of van der Waerden theory is: how large must a set be before every
coloring must contain a monochromatic AP of a prescribed length?

\subsection*{Schur's Equation}

Schur's theorem asks a dual question about additive structure: can we partition the
integers into finitely many sum-free classes, avoiding any monochromatic solution to
$x + y = z$? Schur proved in 1916 that the answer is no for any finite number of
colors, giving rise to the Schur numbers $S(r)$.

\subsection*{Overview}

This paper is organized as follows. Section~\ref{sec:ap} develops the theory of
arithmetic progressions. Section~\ref{sec:vdw-numbers} surveys known van der Waerden
numbers. Section~\ref{sec:proof-vdw} gives the combinatorial proof of van der Waerden's
theorem. Sections~\ref{sec:szemeredi} and~\ref{sec:green-tao} cover the deeper density
and prime theorems. Sections~\ref{sec:schur}--\ref{sec:schur-numbers} treat Schur's
theorem and the Schur numbers. Section~\ref{sec:rado} discusses Rado's generalization.
Section~\ref{sec:connections} surveys connections and open problems.

%% ============================================================
\section{Arithmetic Progressions}\label{sec:ap}
%% ============================================================

\begin{definition}[Arithmetic Progression]
A \emph{$k$-term arithmetic progression} (abbreviated $k$-AP) with first element $a$
and common difference $d > 0$ is the set
\[
  \{a, \, a+d, \, a+2d, \, \ldots, \, a+(k-1)d\} \subset \N.
\]
We call $a$ the \emph{base} and $d$ the \emph{step} of the progression.
A $1$-AP is any singleton; a $2$-AP is any pair $\{a, a+d\}$.
\end{definition}

\begin{definition}[$r$-Coloring and Monochromatic AP]
An \emph{$r$-coloring} of a set $S \subseteq \N$ is any function
$\chi: S \to \{1, 2, \ldots, r\}$.
A $k$-AP $P \subseteq S$ is \emph{monochromatic} under $\chi$ if $\chi$ is constant on $P$.
\end{definition}

The fundamental question is whether every coloring of a sufficiently large initial
segment $\{1, 2, \ldots, N\}$ must produce a monochromatic $k$-AP. Van der Waerden's
theorem answers this affirmatively.

\begin{definition}[Van der Waerden Number]
For integers $k \geq 2$ and $r \geq 2$, the \emph{van der Waerden number} $W(k;r)$
is the smallest positive integer $N$ such that every $r$-coloring of $\{1, \ldots, N\}$
contains a monochromatic $k$-AP.
\end{definition}

The existence of $W(k;r)$ is precisely the content of van der Waerden's theorem.
The notation $W(k; r)$ follows the convention that $k$ is the AP length and $r$ the
number of colors; some sources use $W(r; k)$ or write $W(k_1, k_2, \ldots, k_r)$ for
the asymmetric generalization where color $i$ forbids $k_i$-APs.

\begin{remark}[Trivial cases]
$W(2; r) = r+1$ for all $r \geq 1$: in $\{1, 2, \ldots, r+1\}$, the pigeonhole principle
forces two equal-colored elements, which form a $2$-AP; whereas $\{1,\ldots,r\}$ can be
colored by assigning $i \mapsto i$.
\end{remark}

%% ============================================================
\section{Van der Waerden Numbers}\label{sec:vdw-numbers}
%% ============================================================

Known exact values of $W(k;r)$ are remarkably sparse. Computing them requires both
exhibiting an AP-free coloring of $\{1, \ldots, W-1\}$ (lower bound certificate) and
proving that every coloring of $\{1, \ldots, W\}$ contains a monochromatic AP
(upper bound). Both tasks are computationally expensive and have been tackled primarily
by exhaustive computer search and SAT solvers.

\subsection{Table of Known Values}

\begin{center}
\begin{tabular}{ccc}
\toprule
$k$ & $r$ & $W(k; r)$ \\
\midrule
$3$ & $2$ & $9$ \\
$4$ & $2$ & $35$ \\
$5$ & $2$ & $178$ \\
$6$ & $2$ & $1132$ \\
$3$ & $3$ & $27$ \\
$3$ & $4$ & $293$ \\
$4$ & $3$ & $unknown$ \\
$7$ & $2$ & $unknown$ \\
\bottomrule
\end{tabular}
\end{center}

The value $W(3;2)=9$ has an elegant manual proof: consider any $2$-coloring of
$\{1,\ldots,9\}$. If 1 and 5 receive the same color, then $\{1,5,9\}$ is a
$3$-AP. A careful case analysis confirms that $W(3;2) = 9$.

The value $W(6;2) \geq 1132$ was established by Kouril and Paul (2008) through a
major computer search. Whether $W(6;2) = 1132$ has also been verified
\cite{Kouril2008}, making $W(6;2) = 1132$ exact.

\subsection{Asymptotic Bounds}

Upper bounds on $W(k;2)$ are notoriously hard. The best general result is due to Gowers:

\begin{theorem}[Gowers, 2001 \cite{Gowers2001}]
For all $k \geq 3$,
\[
  W(k; 2) \leq 2^{2^{k+9}}.
\]
\end{theorem}

This tower-of-exponentials bound, while explicit, is far from the truth.
The lower bound of Berlekamp (1968) shows:

\begin{theorem}[Berlekamp, 1968 \cite{Berlekamp1968}]
For every prime $p$, $W(p+1; 2) > p \cdot 2^p$.
\end{theorem}

\begin{proof}[Proof sketch]
Berlekamp constructs an explicit $2$-coloring of $\{1, \ldots, p \cdot 2^p\}$ using
arithmetic in the Galois field $\mathbb{F}_{2^p}$: assign color $0$ to $n$ if the
leading coefficient of a certain polynomial over $\mathbb{F}_2$ is $0$, and color $1$
otherwise. One verifies that this coloring avoids any monochromatic $(p+1)$-AP.
\end{proof}

%% ============================================================
\section{Proof of Van der Waerden's Theorem}\label{sec:proof-vdw}
%% ============================================================

\begin{theorem}[Van der Waerden, 1927 \cite{VdW1927}]\label{thm:vdw}
For every pair of positive integers $k \geq 2$ and $r \geq 2$, there exists a
positive integer $W = W(k; r)$ such that every $r$-coloring of $\{1, \ldots, W\}$
contains a monochromatic $k$-term arithmetic progression.
\end{theorem}

The original proof by van der Waerden (1927) is a purely combinatorial double induction.
We present it in modern language.

\subsection{Combinatorial Proof via Double Induction}

The proof proceeds by induction on $k$ (the AP length), with an inner induction on $r$
(the number of colors) for each fixed $k$.

\begin{definition}[AP-Free Coloring]
A coloring $\chi: \{1,\ldots,N\} \to \{1,\ldots,r\}$ is \emph{$k$-AP-free} if it
contains no monochromatic $k$-AP.
\end{definition}

\begin{lemma}[Base cases]
$W(2; r)$ exists for all $r \geq 1$ (trivially $r+1$ suffices), and $W(k; 1) = k$
for all $k$.
\end{lemma}

\begin{proof}[Proof of Theorem~\ref{thm:vdw}]
We induct on $k$. The base case $k=2$ is trivial.

Assume $W(k-1; r)$ exists for all $r$; we prove $W(k; r)$ exists for all $r$.
Fix $r$ and suppose inductively $W(k; r-1)$ exists; we derive $W(k; r)$.

\medskip
\noindent\textbf{Key construction.}
Set $N_0 = W(k-1; r)$. We will find $W(k; r)$ by a double exponential in $N_0$.

Consider the set $\{1, \ldots, N\}$ with an $r$-coloring $\chi$.
Group it into consecutive blocks of length $N_0$. Within each block,
there is a monochromatic $(k-1)$-AP by the inductive hypothesis.

The crucial observation (the ``density argument'' in van der Waerden's original proof)
is that one can encode the color pattern of $(k-1)$-APs in a block as a ``type''.
There are only finitely many types, so by a second application of the pigeonhole
principle (or an inner Hales--Jewett argument), two blocks with the same type and
matching color pattern can be combined to yield a $k$-AP.

\medskip
Formally, one shows: for each $k$ and $r$,
\[
  W(k; r) \leq W\bigl(k-1;\, r^{W(k;r-1)}\bigr)
\]
by the following argument. Suppose $N \geq W(k-1; r^{W(k;r-1)})$. Given any
$r$-coloring $\chi$ of $\{1,\ldots,N\}$, we can define a derived coloring
$\chi'$ on $\{1,\ldots,N\}$ that encodes the color of runs of length $W(k;r-1)$.
By the inductive hypothesis, $\chi'$ contains a monochromatic $(k-1)$-AP, and
unwrapping the encoding yields a monochromatic $k$-AP under $\chi$.
\end{proof}

\subsection{Connection to Hales--Jewett}

The Hales--Jewett theorem (1963) \cite{HalesJewett1963} provides a powerful generalization:

\begin{theorem}[Hales--Jewett, 1963]\label{thm:hj}
For every $r \geq 2$ and $k \geq 2$, there exists $HJ(k,r)$ such that every
$r$-coloring of the $k$-element $n$-dimensional combinatorial cube
$[k]^n = \{1,\ldots,k\}^n$ (for $n \geq HJ(k,r)$) contains a monochromatic
combinatorial line.
\end{theorem}

Van der Waerden's theorem follows as a corollary of Hales--Jewett by embedding
arithmetic progressions in combinatorial lines. Specifically, every $k$-AP in
$\{1, \ldots, N\}$ corresponds to a combinatorial line in a suitable cube, and
the color of the line is inherited from $\chi$. This gives an alternative (and
in some ways more transparent) proof path.

%% ============================================================
\section{Szemer\'edi's Theorem}\label{sec:szemeredi}
%% ============================================================

Van der Waerden's theorem guarantees monochromatic APs in any coloring of
$\{1,\ldots,N\}$. A much deeper question, posed by Erd\H{o}s and Tur\'an in 1936
\cite{ErdosTuran1936}, asks: if a set $A \subseteq \{1,\ldots,N\}$ has positive
(upper) density, must it contain arbitrarily long APs?

\begin{definition}[Upper Density]
The \emph{upper density} of a set $A \subseteq \N$ is
\[
  \overline{d}(A) = \limsup_{N \to \infty} \frac{|A \cap \{1,\ldots,N\}|}{N}.
\]
\end{definition}

\begin{theorem}[Szemer\'edi, 1975 \cite{Szemeredi1975}]\label{thm:szemeredi}
Every set $A \subseteq \N$ with $\overline{d}(A) > 0$ contains arithmetic progressions
of arbitrary length. That is, for every $k \geq 1$, there exist $a, d$ with $d > 0$
such that $\{a, a+d, \ldots, a+(k-1)d\} \subseteq A$.
\end{theorem}

\begin{remark}[Relation to Erd\H{o}s--Tur\'an]
The \emph{Erd\H{o}s--Tur\'an conjecture} (1936) originally asserted a quantitative
version: a set avoiding $k$-APs must have density going to $0$ at a specific rate,
namely, if $r_k(N)$ denotes the maximum cardinality of a $k$-AP-free subset of
$\{1,\ldots,N\}$, then $r_k(N) = o(N)$.
Szemer\'edi's theorem proves precisely this: $r_k(N)/N \to 0$. The stronger
quantitative refinements (sharp bounds on $r_k(N)$) remain active research,
with Kelley--Meka (2023) establishing $r_3(N) \leq N / (\log N)^{c}$ for some
$c > 0$ for the $3$-AP case.
\end{remark}

\subsection{Earlier Partial Results}

\begin{itemize}
\item \textbf{Roth (1953)}: Proved $r_3(N) = o(N)$, i.e., every dense set contains
  a $3$-AP \cite{Roth1953}.
\item \textbf{Szemer\'edi (1969)}: Extended to $k=4$ using combinatorial methods.
\item \textbf{Szemer\'edi (1975)}: Full result for all $k$ using his Regularity Lemma.
\end{itemize}

\subsection{Polynomial Extensions}

The polynomial Hales--Jewett theorem (Bergelson--Leibman, 1996 \cite{BergelsonLeibman1996})
extends Szemer\'edi's theorem to polynomial progressions:

\begin{theorem}[Bergelson--Leibman, 1996]
Let $p_1, \ldots, p_k \in \Z[t]$ with $p_i(0) = 0$ for all $i$. Then every set of
positive upper density $A \subseteq \N$ contains $a, d$ with $d > 0$ such that
\[
  a + p_1(d),\; a + p_2(d),\; \ldots,\; a + p_k(d) \in A.
\]
\end{theorem}

This subsumes van der Waerden's theorem (linear polynomials $p_i(t) = (i-1)t$) as well
as new configurations such as $\{a, a+d, a+d^2\}$.

%% ============================================================
\section{The Green--Tao Theorem}\label{sec:green-tao}
%% ============================================================

The primes are a sparse set: their density in $\{1,\ldots,N\}$ tends to zero as
$N \to \infty$. Thus Szemer\'edi's theorem does not directly apply to the primes.
Nevertheless, Green and Tao proved:

\begin{theorem}[Green--Tao, 2004 \cite{GreenTao2008}]\label{thm:gt}
The set of prime numbers contains arithmetic progressions of arbitrary length.
\end{theorem}

This theorem was announced in 2004 and published in 2008 after full peer review.
Ben Green and Terence Tao received the Fields Medal (Tao in 2006) in recognition of
this and related work. The proof is one of the most celebrated results in
twenty-first-century mathematics.

\subsection{Outline of the Proof}

The argument proceeds by establishing a \emph{pseudorandomness} property for the primes.
While the primes are sparse, they are ``well distributed enough'' relative to a suitable
majorant. The key steps are:

\begin{enumerate}[label=(\roman*)]
\item \textbf{Pseudorandom majorant.} Using the \emph{Goldston--Y\i ld\i r\i m sieve},
  construct a function $\nu: \N \to [0,\infty)$ that majorizes the indicator of
  primes (up to normalization) and satisfies suitable \emph{linear forms conditions}
  and \emph{correlation conditions}.

\item \textbf{Szemer\'edi-type result for pseudorandom measures.}
  Extend Szemer\'edi's theorem to the setting where the uniform measure is replaced
  by a pseudorandom measure $\nu$: any function $f: \N \to [0,1]$ with
  $\mathbb{E}[f \cdot \nu] \geq \delta > 0$ contains $k$-APs.

\item \textbf{Applying to the primes.}
  The primes satisfy $\mathbb{E}[\mathbf{1}_{\mathrm{primes}} \cdot \nu] \geq \delta$
  for suitable $\nu$ and $\delta$ (by the prime number theorem and the sieve estimates).
  The pseudorandom Szemer\'edi theorem then yields $k$-APs among the primes.
\end{enumerate}

\begin{remark}[Known explicit progressions]
The longest known AP of primes (as of 2024) has 27 terms:
$\{224584605939537911 + k \cdot 81292139 : k = 0, 1, \ldots, 26\}$.
\end{remark}

%% ============================================================
\section{Schur's Theorem}\label{sec:schur}
%% ============================================================

We now turn to Schur's theorem, a result from 1916 predating van der Waerden's theorem
by eleven years. Schur discovered his theorem as a corollary to properties of Fermat's
last theorem modulo primes.

\begin{definition}[Sum-Free Set]
A set $A \subseteq \N$ is \emph{sum-free} if there are no $a, b, c \in A$ (not
necessarily distinct) with $a + b = c$.
\end{definition}

\begin{definition}[Sum-Free Partition]
A partition $\N = A_1 \sqcup A_2 \sqcup \cdots \sqcup A_r$ (restricted to
$\{1, \ldots, n\}$) is \emph{sum-free} (or \emph{Schur partition}) if each class
$A_i$ is sum-free.
\end{definition}

\begin{theorem}[Schur, 1916 \cite{Schur1916}]\label{thm:schur}
For every $r \geq 1$, there exists a positive integer $S$ such that every
$r$-coloring of $\{1, \ldots, S\}$ contains a monochromatic solution to
$x + y = z$ (where $x, y, z \in \{1, \ldots, S\}$ need not be distinct).
\end{theorem}

\begin{proof}
We use Ramsey theory. Let $R = R(3; r)$ be the $r$-color Ramsey number for
triangles (which exists by Ramsey's theorem, 1930). Consider the complete graph
$K_R$ with vertices $\{1, \ldots, R\}$ and color the edge $\{a, b\}$ (with $a < b$)
with color $\chi(b - a)$, where $\chi: \{1,\ldots, R-1\} \to \{1,\ldots,r\}$ is
the given coloring.

By Ramsey's theorem, there exist $1 \leq x < y < z \leq R$ such that the triangle
$\{x, y, z\}$ is monochromatic, i.e., $\chi(y-x) = \chi(z-y) = \chi(z-x)$.
Setting $a = y-x$, $b = z-y$, $c = z-x$, we have $a + b = c$ and all three have
the same color.

Thus we may take $S = R - 1$.
\end{proof}

\begin{remark}
The bound $S \leq R(3; r) - 1$ is far from tight. The actual Schur numbers grow
much more slowly.
\end{remark}

%% ============================================================
\section{Schur Numbers}\label{sec:schur-numbers}
%% ============================================================

\begin{definition}[Schur Number]
The \emph{$r$-th Schur number} $S(r)$ is the largest $n$ such that $\{1, \ldots, n\}$
can be partitioned into $r$ sum-free classes. Equivalently,
\[
  S(r) = \max\bigl\{n : \exists \text{ sum-free $r$-partition of } \{1,\ldots,n\}\bigr\}.
\]
\end{definition}

By Schur's theorem, $S(r)$ is finite for all $r$. The Schur numbers are the
analogue of van der Waerden numbers for the equation $x + y = z$.

\subsection{Known Values}

\begin{center}
\begin{tabular}{ccl}
\toprule
$r$ & $S(r)$ & Status / Year \\
\midrule
$1$ & $1$   & Trivial \\
$2$ & $4$   & Classical \\
$3$ & $13$  & Classical \\
$4$ & $44$  & Baumert 1965 \cite{Baumert1965} \\
$5$ & $160$ & Heule 2017 \cite{Heule2017} (SAT-solver proof) \\
$6$ & $?$   & Open; $536 \leq S(6) \leq 1836$ \\
\bottomrule
\end{tabular}
\end{center}

\subsection{Explicit Partitions}

\begin{example}[$S(1) = 1$]
The single class $\{1\}$ is sum-free since $1 + 1 = 2 \notin \{1\}$.
The set $\{1, 2\}$ cannot be covered: $1 + 1 = 2$ forces a monochromatic solution.
\end{example}

\begin{example}[$S(2) = 4$]
The partition $\{1,4\} \sqcup \{2,3\}$ is sum-free:
$1+1=2\notin\{1,4\}$, $1+4=5\notin\{1,4\}$, $4+4=8\notin\{1,4\}$;
$2+2=4\notin\{2,3\}$, $2+3=5\notin\{2,3\}$, $3+3=6\notin\{2,3\}$.
No $2$-partition of $\{1,2,3,4,5\}$ is sum-free (verified by exhaustion).
\end{example}

\begin{example}[$S(3) = 13$]
An optimal $3$-partition of $\{1,\ldots,13\}$ is:
\[
  K_1 = \{1,4,7,10,13\}, \quad K_2 = \{2,3,11,12\}, \quad K_3 = \{5,6,8,9\}.
\]
Each class is sum-free (the sums $a+b$ for $a,b \in K_i$ always fall outside $K_i$
or exceed $13$). No $3$-partition of $\{1,\ldots,14\}$ is sum-free.
\end{example}

\begin{example}[$S(4) = 44$, Baumert 1965]
Baumert \cite{Baumert1965} established $S(4) = 44$ by computer search, providing
a $4$-partition of $\{1,\ldots,44\}$ and verifying that no $4$-partition of
$\{1,\ldots,45\}$ exists.
\end{example}

\subsection{The Proof of \texorpdfstring{$S(5) = 160$}{S(5) = 160}}

\begin{theorem}[Heule, 2017 \cite{Heule2017}]\label{thm:s5}
$S(5) = 160$.
\end{theorem}

This result was a long-standing open problem. It was resolved by Marijn Heule using
a SAT-solver-assisted computer proof. The proof consists of two parts:

\begin{enumerate}[label=(\roman*)]
\item \textbf{Lower bound}: An explicit $5$-partition of $\{1,\ldots,160\}$ into
  sum-free classes is exhibited. Such a partition can be constructed via the
  \emph{Cube and Conquer} framework (Heule, Kullmann, Manthey 2012
  \cite{HeuleKullmannManthey2012}), combining a greedy seed with backtracking search.

\item \textbf{Upper bound}: The SAT-solver proves that no $5$-partition of
  $\{1,\ldots,161\}$ exists. The SAT instance for $n = 161$, $k = 5$ has
  $161 \times 5 = 805$ Boolean variables $x_{i,c}$ and roughly $\binom{161}{2} \times 5
  \approx 64{,}000$ clauses encoding the Schur constraint $a + b = c \Rightarrow
  \neg(x_{a,j} \wedge x_{b,j} \wedge x_{c,j})$. The solver certifies unsatisfiability
  with a \emph{proof certificate} (DRUP format) verifiable independently.
\end{enumerate}

\begin{remark}[Size of the proof]
The SAT proof for $S(5) = 160$ has a certificate of approximately $200$ terabytes
compressed; it is the largest mathematical proof by size at the time of publication (2017).
\end{remark}

\subsection{The Status of \texorpdfstring{$S(6)$}{S(6)}}

\begin{conjecture}\label{conj:s6}
The exact value of $S(6)$ lies in the interval $[536, 1836]$. The lower bound
$S(6) \geq 536$ is established by an explicit computer-generated $6$-partition of
$\{1, \ldots, 536\}$ (Baumert 1965, refined by later searches). The upper bound
$S(6) \leq 1836$ follows from the general recursive estimate $S(r) \leq r! \cdot e$
(Schur 1916) and SAT-based refinements.
\end{conjecture}

The determination of $S(6)$ is an open problem. The recursive lower bound from the
formula $S(r) \geq 3 \cdot S(r-1) + 1$ gives $S(6) \geq 481$, which is weaker than
the known $536$.

\subsection{Recursive Bounds}

\begin{proposition}[Recursive lower bound]
For all $r \geq 2$,
\[
  S(r) \geq 3 \cdot S(r-1) + 1.
\]
\end{proposition}

\begin{proof}
Given an $r-1$ sum-free partition of $\{1,\ldots,n\}$ (with $n = S(r-1)$), define:
\[
  A_r = \{n+1, n+2, \ldots, 2n+1\}
\]
and rescale the original partition: $A_i' = \{x + 2n + 1 : x \in A_i\}$ for
$i = 1, \ldots, r-1$. The combined partition $(A_1', \ldots, A_{r-1}', A_r)$ gives
an $r$ sum-free partition of $\{1, \ldots, 3n+1\}$:
\begin{itemize}
\item $A_r$ is sum-free: for $a, b \in A_r$, $a + b \geq 2(n+1) > 2n+1$, so $a+b \notin A_r$.
\item Each $A_i'$ inherits sum-freeness from $A_i$ by linearity of the shift.
\end{itemize}
Hence $S(r) \geq 3n + 1 = 3 \cdot S(r-1) + 1$.
\end{proof}

%% ============================================================
\section{Rado's Theorem}\label{sec:rado}
%% ============================================================

Schur's theorem can be viewed as a special case of a much more general result by
Richard Rado (1933) \cite{Rado1933}, which characterizes exactly which systems of
linear equations have the property that every finite coloring of $\N$ produces a
monochromatic solution.

\begin{definition}[Partition Regular]
A matrix $A \in \Z^{m \times n}$ is \emph{partition regular over $\N$} if for every
finite coloring of $\N$, the system $A\mathbf{x} = \mathbf{0}$ has a monochromatic
solution $\mathbf{x} \in \N^n$ (all entries having the same color).
\end{definition}

\begin{definition}[Columns Condition]
A matrix $A = (c_1 \mid c_2 \mid \cdots \mid c_n) \in \Z^{m \times n}$ satisfies
the \emph{columns condition} if there exists a partition $\{1,\ldots,n\} = I_1 \sqcup
I_2 \sqcup \cdots \sqcup I_t$ (into consecutive sets in some ordering) such that:
\begin{itemize}
\item $\sum_{j \in I_1} c_j = \mathbf{0}$,
\item For each $s = 2, \ldots, t$: $\sum_{j \in I_s} c_j \in \mathrm{span}_\Q\{c_j : j \in I_1 \cup \cdots \cup I_{s-1}\}$.
\end{itemize}
\end{definition}

\begin{theorem}[Rado, 1933 \cite{Rado1933}]\label{thm:rado}
A matrix $A \in \Z^{m \times n}$ is partition regular over $\N$ if and only if $A$
satisfies the columns condition.
\end{theorem}

\begin{example}[Schur's equation]
The equation $x + y = z$ corresponds to $A = (1, 1, -1)$. The columns condition
holds: $I_1 = \{3\}$ (so $c_3 = -1$, and $c_1 + c_3 = 0$... let us verify more
carefully). With $I_1 = \{1,2\}$: $c_1 + c_2 = 1 + 1 = 2 \neq 0$. Take
$I_1 = \{3\}$: $c_3 = -1 \neq 0$. Instead: $I_1 = \{1\}$, $I_2 = \{2\}$,
$I_3 = \{3\}$: $c_1 = 1$; $c_1 + c_2 = 2$; check $c_3 = -1 \in \mathrm{span}(c_1, c_2) = \Q$. Yes.
The columns condition is satisfied, confirming partition regularity.
\end{example}

\begin{corollary}
The equation $x + y = z$ is partition regular, recovering Schur's theorem.
The equation $x + y = 3z$ is \emph{not} partition regular (violates columns condition).
\end{corollary}

%% ============================================================
\section{Connections and Open Problems}\label{sec:connections}
%% ============================================================

\subsection{Furstenberg's Ergodic Proof of Szemer\'edi's Theorem}

One of the most conceptually striking developments in this area is Furstenberg's
measure-theoretic proof of Szemer\'edi's theorem (1977) \cite{Furstenberg1977}.

\begin{theorem}[Furstenberg Correspondence Principle]
Let $A \subseteq \N$ with $\overline{d}(A) > 0$. Then there exists a measure-preserving
system $(X, \mathcal{B}, \mu, T)$ and a set $E \in \mathcal{B}$ with $\mu(E) = \overline{d}(A) > 0$
such that for all $k \geq 1$ and all $n \geq 1$:
\[
  d\bigl(A \cap (A-n) \cap (A-2n) \cap \cdots \cap (A-(k-1)n)\bigr) \geq
  \mu(E \cap T^{-n}E \cap T^{-2n}E \cap \cdots \cap T^{-(k-1)n}E).
\]
\end{theorem}

Using the \emph{multiple recurrence theorem}
$\liminf_{N\to\infty} \frac{1}{N}\sum_{n=1}^N \mu(E \cap T^{-n}E \cap \cdots \cap T^{-(k-1)n}E) > 0$,
Furstenberg deduced Szemer\'edi's theorem. This opened the field of \emph{ergodic
Ramsey theory} and inspired all subsequent proofs of Szemer\'edi's theorem.

\subsection{Connections to the Riemann Hypothesis}

A speculative but intriguing connection exists between arithmetic progressions and the
distribution of zeros of the Riemann zeta function $\zeta(s)$.

\begin{conjecture}[Correlation of zeros with APs]\label{conj:rh-ap}
The pair correlation statistics of nontrivial zeros of $\zeta(s)$ (as predicted by the
Montgomery--Odlyzko law and GUE statistics) may encode information about the density of
APs in spectra of certain operators. This connection is currently speculative and not established.
\end{conjecture}

The Green--Tao theorem relies on the prime number theorem (hence on zero-free regions of
$\zeta$), but a direct link between specific properties of Schur or van der Waerden numbers
and the Riemann Hypothesis is not known.

\subsection{Erd\H{o}s--Tur\'an Quantitative Conjecture}

\begin{conjecture}[Erd\H{o}s--Tur\'an, 1936 \cite{ErdosTuran1936}]
If $A \subseteq \N$ is a set such that $\sum_{a \in A} \frac{1}{a}$ diverges, then $A$
contains arithmetic progressions of all lengths.
\end{conjecture}

This conjecture remains open. It would imply the Green--Tao theorem (since $\sum_{p}
\frac{1}{p}$ diverges by Euler), but is strictly stronger. In the direction of the
conjecture, Szemer\'edi's theorem and Green--Tao's theorem give density versions and
prime versions, but the full conjecture is not resolved.

\begin{remark}
Note the distinction: Szemer\'edi's theorem establishes that \emph{positive upper density}
suffices; the Erd\H{o}s--Tur\'an conjecture asks whether \emph{divergence of reciprocal
sum} (a weaker condition than positive density) already suffices. For example, the set
of perfect squares has zero density but divergent harmonic sum is false
($\sum 1/n^2 < \infty$), so the conjecture would not apply there. But a set like
$A = \{n : \exists k, n = 2^k\}$ has divergent sum $\sum_{k\geq 1} 2^{-k} = 1 < \infty$,
so actually converges. A relevant example: primes have zero density but $\sum_p 1/p = \infty$.
\end{remark}

\subsection{Computational Frontiers}

The determination of $W(7;2)$ and $S(6)$ remain major open computational challenges.
Both likely require petabyte-scale SAT solving. The methods of Heule et al.\
(Cube and Conquer, DRUP certificates) continue to be refined.

%% ============================================================
\section*{References}
%% ============================================================

\begin{thebibliography}{99}

\bibitem{VdW1927}
B.L.\ van der Waerden,
\emph{Beweis einer Baudetschen Vermutung},
Nieuw Arch. Wisk.\ \textbf{15} (1927), 212--216.

\bibitem{Schur1916}
I.\ Schur,
\emph{\"{U}ber die Kongruenz $x^m + y^m \equiv z^m \pmod{p}$},
Jahresber. Dtsch. Math.-Ver.\ \textbf{25} (1916), 114--117.

\bibitem{Szemeredi1975}
E.\ Szemer\'edi,
\emph{On sets of integers containing no $k$ elements in arithmetic progression},
Acta Arith.\ \textbf{27} (1975), 199--245.

\bibitem{GreenTao2008}
B.\ Green and T.\ Tao,
\emph{The primes contain arbitrarily long arithmetic progressions},
Ann.\ Math.\ \textbf{167} (2008), 481--547.

\bibitem{Rado1933}
R.\ Rado,
\emph{Studien zur Kombinatorik},
Math.\ Z.\ \textbf{36} (1933), 424--480.

\bibitem{Furstenberg1977}
H.\ Furstenberg,
\emph{Ergodic behavior of diagonal measures and a theorem of Szemer\'edi on arithmetic progressions},
J.\ Anal.\ Math.\ \textbf{31} (1977), 204--256.

\bibitem{Baumert1965}
L.D.\ Baumert,
\emph{A note on Schur's problem},
Unpublished report, Jet Propulsion Laboratory, 1965.

\bibitem{Heule2017}
M.J.H.\ Heule,
\emph{Schur Number Five},
Proc.\ AAAI 2018, arXiv:1711.08076 (2017).

\bibitem{Gowers2001}
W.T.\ Gowers,
\emph{A new proof of Szemer\'edi's theorem},
Geom.\ Funct.\ Anal.\ \textbf{11} (2001), 465--588.

\bibitem{Berlekamp1968}
E.R.\ Berlekamp,
\emph{A construction for partitions which avoid long arithmetic progressions},
Canad.\ Math.\ Bull.\ \textbf{11} (1968), 409--414.

\bibitem{HalesJewett1963}
A.W.\ Hales and R.I.\ Jewett,
\emph{Regularity and positional games},
Trans.\ Amer.\ Math.\ Soc.\ \textbf{106} (1963), 222--229.

\bibitem{Roth1953}
K.F.\ Roth,
\emph{On certain sets of integers},
J.\ London Math.\ Soc.\ \textbf{28} (1953), 104--109.

\bibitem{ErdosTuran1936}
P.\ Erd\H{o}s and P.\ Tur\'an,
\emph{On some sequences of integers},
J.\ London Math.\ Soc.\ \textbf{11} (1936), 261--264.

\bibitem{BergelsonLeibman1996}
V.\ Bergelson and A.\ Leibman,
\emph{Polynomial extensions of van der Waerden's and Szemer\'edi's theorems},
J.\ Amer.\ Math.\ Soc.\ \textbf{9} (1996), 725--753.

\bibitem{HeuleKullmannManthey2012}
M.\ Heule, O.\ Kullmann, and V.\ Manthey,
\emph{Cube and Conquer: Guiding CDCL SAT Solvers by Lookaheads},
Lect.\ Notes Comput.\ Sci.\ \textbf{7155} (2012), 50--65.

\bibitem{Kouril2008}
M.\ Kouril and J.L.\ Paul,
\emph{The van der Waerden number $W(2;6)$ is $1132$},
Experiment.\ Math.\ \textbf{17} (2008), 53--61.

\end{thebibliography}

\end{document}
